\section{High-Arity Workloads (C3, C4)}
\label{sec:high_arity}

Relational facts in real-world domains often exceed the binary subject-predicate-object structure that dominates knowledge graphs. In biomedical databases, a drug--target interaction is annotated with dosage, route, study population, and confidence interval, yielding a 6-ary relation. Sensor networks record timestamped multi-channel readings with provenance metadata, routinely producing tuples of arity 12 or higher. Temporal annotations on facts---valid-from, valid-until, recorded-at, source-id---add further dimensions that push practical arities into double digits.

\paragraph{Reification overhead.}
Standard RDF and property-graph systems encode an $a$-ary fact as $2a+1$ triples via reification or $n$-ary relation patterns. Each of the $a$ role--value pairs requires a separate triple plus the type-assertion triple. Lookup of a single fact therefore requires $O(a)$ edge traversals in the underlying store, and joins across all role--value pairs scale quadratically with arity in pathological cases. This overhead renders high-arity workloads---common in scientific and temporal settings---disproportionately expensive on conventional graph databases.

\paragraph{SheafDB native tuple storage.}
SheafDB eschews reification entirely. The store's fundamental unit is the typed tuple with schema-native arity up to 30. A fact of arity $a$ occupies exactly $a+1$ contiguous cells (arity tag plus $a$ attribute values), yielding $O(1)$ indexed lookup by position within the tuple and $O(1)$ iteration over all attributes. Bulk import of CSV-style records maps directly to internal tuple pages without intermediate triple expansion.

\paragraph{Experimental setup.}
Workloads C3 (point-lookup) and C4 (range-scan) are evaluated at arity 3, 5, and 8, each at scale 1,000 and 10,000 facts. Latency figures are per the micro-benchmark results in Section~\ref{sec:evaluation}.

\paragraph{Complexity analysis.}
SheafDB lookup is $O(1)$ in arity because the complete fact is stored and retrieved as a single section. KG systems pay $O(a)$ due to the $2a+1$ triple expansion. For range-scan workloads, SheafDB's contiguous tuple layout avoids the scatter-read across disjoint triple locations required by reification-based stores. Formal complexity bounds appear in Section~\ref{sec:complexity}.

\endinput
